\chapter{Entrenamiento de las políticas PPO para el robot físico}
\label{cap:ppo-fisico}

Este capítulo describe cómo se formuló, entrenó, seleccionó y ejecutó la capa
correctiva PPO en el robot real entre el 16 de octubre y el 3 de noviembre de 2026. Completa
el contrato de corrección residual (sección~\ref{sec:contrato-residual}) y las
campañas preparatorias en Gazebo y MuJoCo del capítulo de resultados: aquellas usaban un
modelo nominal del actuador, mientras que las de este capítulo se entrenan con
un modelo del actuador identificado con datos del propio robot. La base teórica
de PPO está en la sección de aprendizaje por refuerzo del marco teórico.

\section{Evolución de los entornos de entrenamiento}

La capa correctiva pasó por cuatro entornos, cada uno más cercano al robot:
\begin{enumerate}
    \item \textbf{Banco cinemático en NumPy} (septiembre): observaciones,
    saturaciones y terminaciones sin dinámica ni contacto. Sirvió para validar
    la integración con cinco semillas (11, 23, 37, 53 y 71).
    \item \textbf{Gazebo} (septiembre): la política se conectó al modelo
    dinámico con contacto; ninguna escala de corrección superó a la marcha
    nominal en el criterio conjunto.
    \item \textbf{MuJoCo con modelo nominal} (septiembre y octubre):
    reentrenamiento extendido, con mejora de avance en simulación pero sin
    modelo identificado del actuador.
    \item \textbf{MuJoCo con el actuador real} (desde el 16 de octubre): entorno
    \texttt{NovaRealEnv} (programa \path{entrenar_ppo_real.py}).
    La marcha nominal es la misma que ejecuta el robot, el actuador se
    identifica con sus sensores y la observación se limita a lo que el robot
    puede medir. Es el entorno de todas las campañas de este capítulo.
\end{enumerate}

\section{Entorno de entrenamiento con el actuador real}

\subsection{Modelo físico y temporización}

El simulador es MuJoCo con paso de integración de 2~ms. La política actúa cada
30~ms (15 pasos de integración por paso de control), el mismo periodo que el
bucle de control del robot (33~ms medidos). Desde el 20 de octubre el modelo usa
la geometría medida en el robot (archivo \path{nova_sm3_real.xml}),
Tabla~\ref{tab:geometria-real}. Masas e inercias siguen siendo las del modelo de
referencia NovaSM3.

\begin{table}[H]
\centering
\caption{Geometría medida del robot e incorporada al modelo de entrenamiento.}
\label{tab:geometria-real}
\small
\begin{tabular}{lrr}
\toprule
Magnitud & Robot real (mm) & Modelo anterior (mm) \\
\midrule
Fémur (eje de cadera a eje de rodilla) & 108 & 107,95 \\
Tibia hasta la punta del pie & 155 & 153,89 \\
Desplazamiento lateral de la coxa & 35 & 38,1 \\
Separación de caderas delante--detrás & 190 & 180 \\
Separación de caderas izquierda--derecha & 77 & 120 \\
Largo del cuerpo & 265 & 230 \\
\bottomrule
\end{tabular}
\end{table}

\subsection{Marchas nominales}

La marcha nominal $\mathbf{q}^{\text{nom}}(\phi)$ de la ecuación~\eqref{eq:residual}
es una de las cinco de la Tabla~\ref{tab:marchas-nominales}:
\begin{itemize}
    \item \textbf{Gateo}: una pata en el aire cada vez, en el orden delantera
    izquierda, trasera derecha, delantera derecha, trasera izquierda, con una
    pausa entre fases.
    \item \textbf{Paso}: igual secuencia pero continua. Se dejó el 20 de octubre
    porque, con las coxas fijas, resultaba un gateo continuo visualmente
    indistinguible del gateo; se retomó el 22 de octubre (v16) con el actuador
    medido.
    \item \textbf{Trote}: patas diagonales a la vez, dos en el aire y dos apoyadas.
    \item \textbf{Carrera} (desde el 27 de octubre): también por pares diagonales,
    con un ciclo más corto y una fracción de apoyo de 0,45, de modo que en la
    referencia hay un breve instante sin pies apoyados.
    \item \textbf{Galope transverso} (desde el 27 de octubre): las patas de cada eje
    casi a la vez, con fases trasera izquierda 0, trasera derecha 0,1, delantera
    derecha 0,5 y delantera izquierda 0,6, y fracción de apoyo de 0,40. Es la única
    marcha que mueve coxas: las traseras se abren $4^\circ$ al levantar, y las
    delanteras quedan fijas, porque con la coxa 2 en movimiento la pata cedía y
    con una sola delantera moviéndose el galope retrocedía en simulación.
\end{itemize}

\begin{table}[H]
\centering
\caption{Parámetros de las marchas nominales actuales.}
\label{tab:marchas-nominales}
\small
\begin{tabular}{lrrrrr}
\toprule
Marcha & Ciclo (s) & Fémur & Tibia & Apertura & Peso de avance $w_{\text{av}}$ \\
\midrule
Gateo & 2,64 & $\pm15^\circ$ & $10^\circ$ & $8^\circ$ & 40 \\
Paso & 2,0 & $\pm15^\circ$ & $10^\circ$ & $8^\circ$ & 20 \\
Trote & 1,0 & $\pm12^\circ$ & $12^\circ$ & $5^\circ$ & 10 \\
Carrera & 0,8 & $\pm12^\circ$ & $12^\circ$ & $5^\circ$ & 10 \\
Galope & 0,8 & $\pm18^\circ$ & $7^\circ$ & $5^\circ$ & 30 \\
\bottomrule
\end{tabular}
\end{table}

Desde el 20 de octubre las marchas usan una \emph{postura abierta}: el punto
neutro de los fémures delanteros se adelanta y el de los traseros se atrasa según la
columna de apertura. Con la geometría real, las patas de
un mismo lado se acercaban hasta 69~mm bajo el centro del cuerpo y llegaban a
chocar; con la apertura el hueco mínimo entre pies sube a 133~mm.

Desde el 27 de octubre la tibia de cada pata amplía su movimiento con una ganancia
propia (tibias 1 a 4: 2,0; 1,4; 1,6 y 1,8), y las tibias 2 y 4 se mueven en sentido
de servo contrario a las 1 y 3. Es la postura de trabajo aprobada por el operador. Estas
ganancias se incorporaron a la simulación desde la campaña v24; la lectura de los
sensores se divide por la misma ganancia, de modo que la política ve el ángulo del
modelo. Con ellas, el galope anterior (fémur $16^\circ$, tibia $14^\circ$) se caía
siempre, y su marcha base se reajustó a los valores de la tabla
(sección~\ref{sec:viabilidad-resultados}).

\subsection{Acción}

La acción es un vector de 12 componentes, una por articulación. Las cuatro
coxas están enmascaradas (corrección cero), por lo que la política corrige
ocho articulaciones: los cuatro fémures y las cuatro tibias. La salida de la
red se recorta a $[-1,1]$ y se escala a $\pm0{,}08$~rad ($\pm4{,}6^\circ$); la
corrección aplicada no puede cambiar más de $0{,}02$~rad ($1{,}15^\circ$) por
paso respecto a la anterior. Además, como el firmware del robot, la orden total
no puede variar más de $3^\circ$ por paso y se recorta a los topes medidos de
cada servo.

\subsection{Observación}

La observación incluye solo magnitudes medibles en el robot suspendido, sin IMU. Desde la versión 6 tiene 22 componentes:
\begin{itemize}
    \item desviación de las ocho articulaciones de fémur y tibia respecto a su
    postura de referencia (8), con ruido gaussiano de $1^\circ$ y un sesgo
    aleatorio de hasta $\pm1^\circ$ por episodio, como los potenciómetros;
    \item fase de la marcha codificada como $\sin(2\pi\phi)$ y $\cos(2\pi\phi)$ (2);
    \item corrección del paso anterior normalizada por $0{,}08$~rad (12).
\end{itemize}
Las versiones 2 a 5 observaban solo tres articulaciones (fémures izquierdos y
tibia delantera izquierda), las únicas con sensor fiable en esa fecha, con 17
componentes.

\subsection{Recompensa}

La recompensa por paso es
\begin{multline}
    r_t = 1 + w_{\text{av}}\,\Delta x - w_{\text{act}}\,(\varphi^2+\vartheta^2)
    - 25\,(h-0{,}235)^2 - 0{,}25\,\lVert\mathbf{e}_q\rVert^2 - 0{,}04\,\lVert\Delta\mathbf{q}\rVert^2 \\
    - w_{\text{lat}}\,y^2 - w_{\text{giro}}\,\psi^2
    - w_{\text{suav}}\sum_j\Big(\frac{\Delta\mathbf{q}_{t,j}-\Delta\mathbf{q}_{t-1,j}}{0{,}02}\Big)^2
    - w_{\text{cox}}\,\overline{e_{\text{coxa}}^2},
    \label{eq:recompensa-real}
\end{multline}
donde $\Delta x$ es el avance del cuerpo en el paso (m); $\varphi$ y $\vartheta$,
el alabeo y el cabeceo (rad); $h$, la altura del cuerpo (m); $\mathbf{e}_q$, el
error articular respecto a la marcha nominal; $\Delta\mathbf{q}$, la corrección;
$y$ y $\psi$, el desvío lateral (m) y el giro (rad) acumulados desde el inicio;
y $\overline{e_{\text{coxa}}^2}$, la media del cuadrado de la desviación de las
coxas respecto a su orden (en grados). Además, al completar cada ciclo se resta
$w_{\text{gc}}\,(\Delta\psi_{\text{ciclo}})^2$, con el giro del ciclo en grados,
porque sin IMU la política no observa el rumbo y tendía a curvarse en caminatas
largas. Los pesos de cada versión están en las
tablas~\ref{tab:campanas-ppo} y~\ref{tab:campanas-ppo-2}.

El término de suavidad (versión 11 en adelante) se añadió porque en el robot la
primera PPO del paso daba correcciones más rápidas de lo que los servos podían
seguir; el término de coxas, porque esa misma política sacudía el cuerpo y las
coxas, fijas en la simulación pero flexibles en el robot, se desviaban el doble
que con la marcha nominal.

\subsection{Terminación y duración de los episodios}

Un episodio termina si la altura del cuerpo sale de $[0{,}12;\,0{,}40]$~m, si
el alabeo o el cabeceo superan $0{,}6$~rad o si la observación deja de ser
finita. Si no, se trunca al completar 6 ciclos de marcha (12 en el trote, por
su ciclo más corto). Cada episodio comienza con una rampa desde la postura de
reposo hasta la postura de la marcha, durante la cual la política no actúa.

\subsection{Modelo del actuador y aleatorización}

Entre la orden y la simulación se interpone el modelo del actuador
identificado en el robot: un retardo puro (cola de 3 pasos, $\approx$99~ms), una
ganancia por articulación, una zona muerta ($\approx0{,}4^\circ$) y el límite de
$3^\circ$ por paso. Se identificó por primera vez el 16 de octubre con las
caminatas nominales registradas por los sensores del robot (retardo de 65 a
100~ms, ganancia de 0,88 a 1,0, histéresis de 0,6 a $1{,}5^\circ$) y se
actualizó después con escalones, senoides y barridos de rango.

En cada episodio se aleatorizan el retardo (2 a 4 pasos; desde el 21 de
octubre, el intervalo medido en el robot, que en las políticas elegidas es de 1 a 4
pasos, es decir, de 30 a 120~ms), la ganancia de cada servo
($\pm5$~\%) y el sesgo de los sensores ($\pm1^\circ$), según la estrategia de
aleatorización del dominio descrita en el marco teórico.

\section{Red neuronal e hiperparámetros}

Se usó la implementación de PPO de Stable-Baselines3~2.3.2~\cite{raffin2021sb3}
sobre PyTorch~2.3.1, con la política \texttt{MlpPolicy}: actor y crítico son
redes independientes, cada una con dos capas ocultas de 64 neuronas y activación
$\tanh$. El actor da la media de la gaussiana; las desviaciones son parámetros
aprendidos independientes del estado. La Tabla~\ref{tab:hiperparametros-ppo}
resume los hiperparámetros de las campañas desde la versión 10; en las versiones
2 a 6 no se registró la configuración de PPO en los metadatos.

\begin{table}[H]
\centering
\caption{Hiperparámetros de PPO en las campañas con el actuador real (versión 10 en adelante).}
\label{tab:hiperparametros-ppo}
\small
\begin{tabular}{lll}
\toprule
Hiperparámetro & Valor & Significado \\
\midrule
Tasa de aprendizaje & $1\cdot10^{-4}$ & paso del optimizador Adam \\
Pasos por recolección $N$ & 2048 & datos por actualización \\
Tamaño de minilote & 256 & muestras por paso de gradiente \\
Épocas por actualización $K$ & 10 & pasadas sobre cada recolección \\
Factor de descuento $\gamma$ & 0,99 & horizonte de $\approx$100 pasos (3~s) \\
Parámetro GAE $\lambda$ & 0,95 & compromiso sesgo--varianza \\
Recorte $\epsilon$ & 0,2 & cambio máximo de la razón $\rho_t$ \\
Coeficiente del crítico $c_1$ & 0,5 & peso del error de valor \\
Coeficiente de entropía $c_2$ & 0 & sin bonificación de exploración \\
Norma máxima del gradiente & 0,5 & recorte del gradiente \\
KL objetivo & 0,03 & corta las épocas si la política cambia mucho \\
Copias intermedias & cada $10^5$ pasos & para la selección; más densas en campañas cortas \\
Evaluación final & semillas 801--805 & 5 episodios deterministas \\
\bottomrule
\end{tabular}
\end{table}

\section{Campañas de entrenamiento y semillas}

La Tabla~\ref{tab:campanas-ppo} recoge las primeras campañas con el actuador real.
Cada fila es una configuración entrenada con varias semillas independientes; la
columna de pasos indica el máximo alcanzado (en las continuaciones, contando los
pasos de la política de partida). Las semillas 11, 23, 37, 53 y 71 son las
congeladas en el protocolo; las demás se añadieron para ampliar la muestra. Las
tablas~\ref{tab:campanas-ppo} y~\ref{tab:campanas-ppo-2} dividen las campañas en
dos periodos. Las continuaciones parten de la mejor política de una campaña anterior y siguen
entrenando con la configuración nueva.

\begin{table}[H]
\centering
\caption{Campañas PPO con el actuador real, 16 a 22 de octubre de 2026.}
\label{tab:campanas-ppo}
\scriptsize
\setlength{\tabcolsep}{3pt}
\begin{tabular}{p{0.1\textwidth}p{0.06\textwidth}p{0.07\textwidth}p{0.19\textwidth}p{0.08\textwidth}p{0.4\textwidth}}
\toprule
Versión & Fecha & Marcha & Semillas & Pasos & Cambio principal y resultado \\
\midrule
v2, v3 & 16/10 & gateo & 11, 23, 37, 53, 71 & 0,15 M & Primer entorno con el actuador real; 3 articulaciones observadas. Todas superan el avance nominal en simulación. \\
v4 & 16/10 & gateo & 11, 23, 37, 53, 71 & 0,2 M & Recompensa de estabilidad ($w_{\text{av}}=10$, $w_{\text{act}}=30$, $w_{\text{lat}}=20$, $w_{\text{giro}}=10$). Semilla 37: $-54$~\% de deriva y $-31$~\% de inclinación en simulación; primera ejecutada en el robot. \\
v5 & 16/10 & gateo & 11, 23, 37, 53, 71 & 0,2 M & Actuador v4 (barrido de rango). No mejora la base. \\
v6 & 19/10 & gateo & 11, 23, 37, 53, 71 & 0,2 M & Observación de las 8 articulaciones (22 componentes). \\
v10 & 19/10 & gateo, paso & 11, 23, 37, 53, 71 & 3,0--3,2 M & Aleatorización del actuador, $w_{\text{gc}}=5$, tasa $10^{-4}$. Elegidas: gateo semilla 11 (0,3 M) y paso semilla 71 (2,3 M). \\
v11 & 20/10 & paso & 11, 23, 37, 53, 61, 71, 83, 97 & 5,3 M & Suavidad ($w_{\text{suav}}=0{,}3$) y coxas ($w_{\text{cox}}=1$). Mejora a la v10, pero no a la base en el robot. \\
v12 & 20/10 & gateo & 11, 37, 41, 59 & 4,8 M & Geometría real. Elegida: semilla 41 (0,5 M). \\
trote & 20/10 & trote & 11, 23, 37, 53, 61, 71, 83, 97 & 3,0 M & Primera PPO del trote, recompensa de la v11. \\
v13 & 20/10 & gateo & 73, 89, 107, 127, 149, 163, 181, 199, 223, 251 & 3,0 M & Postura abierta $8^\circ$ y tibia $14^\circ$; pesos por defecto ($w_{\text{av}}=18$, $w_{\text{act}}=8$, sin $w_{\text{lat}}$ ni $w_{\text{giro}}$). Evaluadas solo en simulación. \\
v14 & 21/10 & gateo, trote & gateo 13, 29, 41, 43, 67; trote 19, 31, 37, 47, 89 & 3,0 M & Postura abierta; continúan copias de 0,9 M. Elegida gateo: semilla 43 (1,7 M). Ninguna de trote supera la base. \\
v15 & 21/10 & trote & 19, 211, 223, 241, 263, 281, 307 & 3--5 M & $w_{\text{av}}=20$, $w_{\text{suav}}=0{,}1$. Elegida: semilla 307 (0,1 M). \\
ciclo 1 & 21/10 & gateo, trote & 401$^{\dagger}$ & +1 M & Continúan las ganadoras en el robot. Elegida gateo: copia de 2,7 M. \\
ciclo 2 & 21/10 & gateo, trote & 401$^{\dagger}$ & +1 M & Actuador medido en el robot en la vuelta anterior. Elegida gateo: copia de 2,4 M. \\
v16 & 22/10 & paso & continuación: 503, 509, 521, 523; desde cero: 541, 547, 557, 563 & 6,4 M / 3,0 M & Retoma el paso con el actuador medido (ciclo 2) y la postura corregida ($8^\circ$); las continuaciones parten de la v11 semilla 71 (4,9 M). Elegida: semilla 563 (1,6 M); en simulación, deriva de 0,021 m frente a 0,391 m de la nominal y avance $-12$~\%. \\
\bottomrule
\multicolumn{6}{p{0.95\textwidth}}{\footnotesize $^{\dagger}$ Se lanzaron con las semillas 401, 409, 419 y 421, pero las cuatro ejecuciones resultaron idénticas: al continuar desde una política guardada, la biblioteca restaura la semilla almacenada en ella e ignora la indicada. Cada una de estas campañas equivale a una sola ejecución. El error se corrigió el 22 de octubre de 2026 (se impone la semilla pedida tras cargar la política).}
\end{tabular}
\end{table}

Se observaron tres patrones. Primero, la variabilidad entre semillas fue grande:
en la campaña v15 del trote, la recompensa media final de las siete semillas
osciló entre $-60$ y $+64$ con la misma configuración, lo que confirma la
necesidad de entrenar varias semillas. Segundo, la mejor copia de una campaña no
siempre es la última: varias de las elegidas tienen entre $0{,}1$ y $0{,}5$~M de
pasos, porque al seguir entrenando la política ganaba avance a costa de
inclinación o de suavidad. Tercero, las continuaciones heredan lo aprendido: las
semillas que partían de una política previa alcanzaban en pocos cientos de miles
de pasos una recompensa que las semillas nuevas tardaban millones de pasos en
lograr. La Figura~\ref{fig:curvas-semillas} muestra estos patrones en tres
campañas.

\begin{figure}[H]
\centering
\includegraphics[width=\textwidth]{Resultados/ppo_curvas_semillas.png}
\caption{Recompensa media por episodio durante el entrenamiento, una línea por
semilla, en tres campañas: gateo v14 (continuación desde copias de 0,9~M pasos),
trote v15 (la semilla 19 continúa una política previa hasta 5~M pasos) y primer
trote (8 semillas desde cero). Cada panel tiene su propia escala vertical, porque
los pesos de la recompensa difieren entre campañas. Curvas suavizadas con una
media móvil de 15 puntos.}
\label{fig:curvas-semillas}
\end{figure}

Del 26 de octubre en adelante, las campañas se dedicaron a añadir la carrera y el
galope, a incorporar la postura de trabajo nueva y a elegir las políticas por
viabilidad (sección de selección de copias), Tabla~\ref{tab:campanas-ppo-2}. Los
pesos de la recompensa son los de la v4 con suavidad 0,3 y coxas 1, cambiando el
peso de avance $w_{\text{av}}$ y, en el galope, $w_{\text{lat}}=30$.

\begin{table}[H]
\centering
\caption{Campañas PPO con el actuador real, 26 de octubre a 3 de noviembre de 2026.}
\label{tab:campanas-ppo-2}
\scriptsize
\setlength{\tabcolsep}{3pt}
\begin{tabular}{p{0.07\textwidth}p{0.07\textwidth}p{0.16\textwidth}p{0.11\textwidth}p{0.49\textwidth}}
\toprule
Versión & Fecha & Marcha & Semillas & Cambio principal y resultado \\
\midrule
v17 & 26/10 & paso & 4 & Continúan la v16 semilla 563; ninguna la supera (mejor puntuación 2,50 frente a 2,45). \\
v18 & 27/10 & carrera, galope & 4 y 4 & Primeras campañas de ambas marchas, desde cero. \\
v19--v21 & 27/10 & galope & 4 cada una & Pruebas de coxas en movimiento. Mejor: v21 semilla 827 (0,1~M), 90~\% de viabilidad. \\
v22 & 27/10 & gateo, paso, trote, carrera & 4, 6, 4, 4 & Primera selección por viabilidad; 100~\% en las cuatro, aún sin las ganancias de tibia en la simulación. Paso con $w_{\text{av}}=20$. \\
v23 & 27/10 & galope & 8 & Detenida antes de su política final; no se usó. \\
v24 & 27--28/10 & gateo, paso, trote, carrera, galope & 4 cada una & Ganancias de tibia en la simulación. 100~\%: paso semilla 4109, trote 4209 y carrera 4309, las tres a 0,2~M. Gateo $\leq$55~\%. La marcha base del galope se cae en todos los episodios. \\
v25 & 28/10 & gateo, galope & 4 y 4 & $w_{\text{av}}=25$ y 15; galope con fémur $18^\circ$ y tibia $7^\circ$. Galope 65~\%, gateo $\leq$55~\%. \\
v26 & 28/10 & gateo, galope & 4 y 4 & $w_{\text{av}}=40$ y 30. Galope 85~\%, gateo $\leq$55~\%. \\
v27 & 28/10 & gateo, galope & 4 y 6 & Elegida gateo: semilla 7013 (4,1~M), 100~\%; continúa la v22. Galope 90~\%. \\
v28 & 28/10 & galope & 8 & Continúan copias de la v27. Elegida: semilla 8417 (0,175~M), 90~\%. \\
v29, v30 & 29/10 & trote; las cinco & 3; 1 cada una & Solo para registrar pérdidas y recompensa por episodio, 1~M de pasos. Políticas no usadas. \\
v31 & 3/11 & las cinco & 1 cada una & Solo para medir la convergencia (sección~\ref{sec:convergencia-marchas}); 3~M de pasos, galope 5~M. Políticas no usadas. \\
\bottomrule
\end{tabular}
\end{table}

\section{Convergencia del entrenamiento por marcha}
\label{sec:convergencia-marchas}

Para medir cuándo deja de mejorar cada marcha se realizó el 3 de noviembre de 2026
una campaña específica (v31): las cinco marchas se entrenaron desde cero en MuJoCo,
con los hiperparámetros de la Tabla~\ref{tab:hiperparametros-ppo}, durante 3~M de pasos
(5~M en el galope, que a 1~M aún no se había estabilizado). Se usó una sola semilla por
marcha. Las políticas de esta campaña no se probaron en el robot: la campaña se
hizo únicamente para observar la convergencia. Además de la recompensa total, se
registró en cada episodio el valor de cada término de la
ecuación~\eqref{eq:recompensa-real}.

\subsection{Composición de la recompensa}

La recompensa tiene dos premios y nueve castigos. Los premios corresponden a lo que se
quiere aumentar: seguir de pie (el término constante $+1$ por paso, que se pierde al
caer porque el episodio termina) y avanzar. Los castigos corresponden a lo que se quiere
llevar a cero: inclinación, variación de altura, alejamiento de la marcha nominal, tamaño
de la corrección, desvío lateral, giro, giro por ciclo, brusquedad y desviación de las
coxas. Los pesos se ajustaron a mano entre campañas a partir de los problemas observados
en la simulación y en el robot; no proceden de una deducción analítica.

Escribir un objetivo como premio o como castigo es equivalente para el aprendizaje,
porque la política se ajusta según la diferencia entre acciones y no según el signo
absoluto de la recompensa. Por ello, el valor negativo de la recompensa total no indica
por sí solo un mal comportamiento: lo que se interpreta es su evolución. La
Tabla~\ref{tab:terminos-recompensa} muestra que, al final del entrenamiento, los castigos
siguen superando a los premios en todas las marchas.

\begin{table}[H]
\centering
\caption{Principales términos de la recompensa por episodio al final del entrenamiento (media del último 10~\% de los episodios, en puntos).}
\label{tab:terminos-recompensa}
\small
\begin{tabular}{lrrrrr}
\toprule
Marcha & Seguir de pie & Avance & Brusquedad & Coxas & Giro por ciclo \\
\midrule
Gateo & $+561$ & $+24$ & $-239$ & $-383$ & $-9$ \\
Paso & $+435$ & $+10$ & $-234$ & $-283$ & $-15$ \\
Carrera & $+357$ & $+3$ & $-125$ & $-313$ & $-10$ \\
Trote & $+429$ & $+6$ & $-461$ & $-315$ & $-13$ \\
Galope & $+357$ & $+5$ & $-116$ & $-729$ & $-114$ \\
\bottomrule
\end{tabular}
\end{table}

\subsection{Lectura de las gráficas}

Las Figuras~\ref{fig:convergencia-gateo} a~\ref{fig:convergencia-galope} tienen los mismos
seis paneles. Los dos primeros muestran la recompensa total frente a los pasos de
entrenamiento y frente a las épocas; los dos siguientes, la suma de los castigos con
los mismos ejes, en valor positivo, de modo que en ellos una curva más baja es mejor. Un
paso es una decisión de la política (cada 30~ms) y una época es una pasada de
optimización sobre los 2048 pasos de una recolección. El quinto panel relaciona ambas
escalas: es una recta porque cada recolección usó sus 10 épocas y el umbral de divergencia
KL no interrumpió ninguna, de modo que 1~M de pasos equivale a unas 4880 épocas. El sexto
panel separa los castigos de mayor peso y permite ver cuáles reduce la política y cuáles
apenas cambian.

Se adoptó el siguiente criterio de convergencia: sobre la media móvil de 100
episodios, es el primer punto a partir del cual la curva ya no sale de una banda de
$\pm2$~\% de su recorrido total alrededor de su valor final (media del último 10~\% de
los episodios). Una curva cuyo punto cayera en el último 20~\% del entrenamiento se
consideraría no convergida; es el caso del galope. Una primera versión del criterio usaba una
banda de $\pm5$~\%, pero situaba el punto cuando la curva aún subía de forma visible: como
la banda se mide sobre todo el recorrido, que empieza muy abajo, quedaba demasiado ancha.
El criterio es una convención de este trabajo, y los puntos obtenidos corresponden a una sola semilla.

\begin{table}[H]
\centering
\caption{Punto de convergencia de la recompensa y evolución de la recompensa, el castigo y el ruido por episodio (desviación típica del primer y del último 10~\% de los episodios; recompensa y castigo inicial: primer 1~\%).}
\label{tab:convergencia-marchas}
\small
\begin{tabular}{lrrrrr}
\toprule
Marcha & Pasos (M) & Épocas & Recompensa & Castigo & Ruido \\
\midrule
Gateo & 0,89 & 4353 & $-1438 \to -103$ & $2028 \to 688$ & $242 \to 16$ \\
Paso & 0,95 & 4641 & $-1214 \to -127$ & $1663 \to 572$ & $182 \to 18$ \\
Carrera & 1,23 & 5989 & $-1740 \to -215$ & $2103 \to 574$ & $365 \to 14$ \\
Trote & 2,06 & 10\,030 & $-1750 \to -439$ & $2186 \to 874$ & $321 \to 16$ \\
Galope & no converge & --- & $-4699 \to -742$ & $5072 \to 1104$ & $1281 \to 281$ \\
\bottomrule
\end{tabular}
\end{table}

\subsection{Resultados}

La Tabla~\ref{tab:convergencia-marchas} resume los puntos de convergencia. La recompensa
y el castigo convergen en el mismo punto o muy cerca (en el gateo, el castigo a 0,80~M), porque los premios apenas varían
durante el entrenamiento: la recompensa sube, sobre todo, porque bajan los castigos. El
gateo y el paso convergen antes de 1~M de pasos, la carrera hacia 1,2~M y el trote hacia
2,1~M. El galope no converge en 5~M de pasos: aprende en dos etapas, con una meseta entre
1 y 2~M de pasos y un salto hacia los 2~M, tras el cual sigue mejorando lentamente hasta
casi el final (su curva solo queda dentro de la banda a partir de 4,8~M), y es la marcha que conserva más ruido al final; también es la única con caídas durante el
entrenamiento (7 en 14\,010 episodios).

Por término, el giro por ciclo y la brusquedad son los castigos que más disminuyen. El
castigo de las coxas baja menos en las cuatro primeras marchas, y en el galope prácticamente no
cambia. La política no comanda las coxas: solo influye en ellas de forma indirecta,
a través de las sacudidas del cuerpo.

La convergencia indica que la recompensa dejó de mejorar, no que la marcha sea
mejor para el objetivo de avanzar. El avance por episodio \emph{disminuyó} durante el
entrenamiento en todas las marchas: de 0,74 a 0,60~m en el gateo, de 0,67 a 0,51~m en el
paso, de 0,73 a 0,62~m en el trote, de 0,58 a 0,27~m en la carrera y de 0,57 a 0,15~m en el
galope (primer 1~\% frente a último 10~\% de los episodios). El peso del avance es pequeño frente
al de los castigos, de modo que la política aprende a moverse con más suavidad y con
menos desvío a costa de avanzar menos. Este resultado es coherente con el hecho de que la mejor copia
de una campaña no siempre sea la última, y justifica la selección que se describe a
continuación.

\begin{figure}[p]
\centering
\includegraphics[width=\textwidth,height=0.86\textheight,keepaspectratio]{Resultados/ppo_convergencia_gateo.png}
\caption{Convergencia del entrenamiento. Gateo (semilla 4009, 3~M de pasos). Tono claro: valor de cada episodio; tono oscuro: media móvil de 100 episodios; línea discontinua: punto de convergencia. Ejes verticales recortados al percentil 0,5--99,5.}
\label{fig:convergencia-gateo}
\end{figure}

\begin{figure}[p]
\centering
\includegraphics[width=\textwidth,height=0.86\textheight,keepaspectratio]{Resultados/ppo_convergencia_paso.png}
\caption{Convergencia del entrenamiento. Paso (semilla 4109, 3~M de pasos). Tono claro: valor de cada episodio; tono oscuro: media móvil de 100 episodios; línea discontinua: punto de convergencia. Ejes verticales recortados al percentil 0,5--99,5.}
\label{fig:convergencia-paso}
\end{figure}

\begin{figure}[p]
\centering
\includegraphics[width=\textwidth,height=0.86\textheight,keepaspectratio]{Resultados/ppo_convergencia_carrera.png}
\caption{Convergencia del entrenamiento. Carrera (semilla 4309, 3~M de pasos). Tono claro: valor de cada episodio; tono oscuro: media móvil de 100 episodios; línea discontinua: punto de convergencia. Ejes verticales recortados al percentil 0,5--99,5.}
\label{fig:convergencia-carrera}
\end{figure}

\begin{figure}[p]
\centering
\includegraphics[width=\textwidth,height=0.86\textheight,keepaspectratio]{Resultados/ppo_convergencia_trote.png}
\caption{Convergencia del entrenamiento. Trote (semilla 4207, 3~M de pasos). Tono claro: valor de cada episodio; tono oscuro: media móvil de 100 episodios; línea discontinua: punto de convergencia. Ejes verticales recortados al percentil 0,5--99,5.}
\label{fig:convergencia-trote}
\end{figure}

\begin{figure}[p]
\centering
\includegraphics[width=\textwidth,height=0.86\textheight,keepaspectratio]{Resultados/ppo_convergencia_galope.png}
\caption{Convergencia del entrenamiento. Galope (semilla 6403, 5~M de pasos). Tono claro: valor de cada episodio; tono oscuro: media móvil de 100 episodios. Sin línea discontinua, porque no converge con el criterio usado. Ejes verticales recortados al percentil 0,5--99,5.}
\label{fig:convergencia-galope}
\end{figure}

\section{Selección de copias}

Cada campaña deja una copia de la política cada $10^5$ pasos. La recompensa de
entrenamiento no basta para elegir entre ellas, porque los episodios son cortos
(6 ciclos) y sus valores dependen de los pesos de cada versión. Por eso todas
las copias se evalúan con el script \path{seleccionar_copias.py} en caminatas
simuladas largas de 60 ciclos ($\approx$3~min), con dos semillas de evaluación
(901 y 902) y en las mismas condiciones que la marcha nominal. Una copia es
\emph{válida} si avanza más que la marcha nominal, no se desvía más del límite
fijado y no se cae. Entre las válidas se elige la de menor puntuación
\begin{equation}
    P = \frac{d_y}{0{,}1\,\text{m}} + \varphi_{\text{rms}} + \vartheta_{\text{rms}}
    - \frac{x}{x_{\text{nom}}} + 2\left(\frac{c_{\text{rms}}}{c_{\text{rms,nom}}}-1\right) + S,
    \label{eq:puntuacion-seleccion}
\end{equation}
con $d_y$ el desvío lateral final, $\varphi_{\text{rms}}$ y $\vartheta_{\text{rms}}$
el alabeo y el cabeceo RMS en grados, $x$ el avance, $c_{\text{rms}}$ la
desviación RMS de las coxas y $S$ el cambio cuadrático medio de la corrección
(suavidad); el subíndice nom indica la marcha nominal. La copia ganadora se
guarda en \path{ppo_v10_robot/} junto con sus metadatos (marcha, amplitudes,
apertura, versión del actuador y resultado de la evaluación), que el programa
del robot lee para reproducir exactamente la misma configuración.

Desde el 27 de octubre, esta selección se completa con un criterio más directo, la
\emph{viabilidad} (\path{evaluar_viabilidad.py}), que mide si la copia puede
sustituir a la marcha nominal cuando cambian las condiciones. Sobre 20 episodios con
semillas 1001 a 1020, aleatorización del dominio y del actuador, un episodio $i$ es
viable si cumple a la vez
\begin{equation}
    \text{sin caída},\quad
    x_i \geq 0{,}9\,x_{\text{nom}},\quad
    \left(d_i \leq d_{\text{nom}} \ \text{o}\ d_i < 0{,}5~\text{m}\right),\quad
    \theta_i < 15^\circ,
    \label{eq:viabilidad}
\end{equation}
con $x_i$ el avance, $d_i$ la deriva lateral y $\theta_i$ la inclinación máxima del
episodio. La viabilidad es el porcentaje de episodios que lo cumplen. Sus resultados
están en la sección~\ref{sec:viabilidad-resultados}.

\section{Ejecución en el robot}

\subsection{Programa de control y realimentación}

En el robot, la política se ejecuta con el programa de control
\path{ejecutar_politica_robot.py} (en el PC durante los ensayos de octubre). En cada paso de 30~ms el programa calcula la
referencia nominal, construye la observación con los sensores, evalúa la red
(media de la gaussiana, sin muestreo), aplica las saturaciones, convierte los
ángulos del modelo a grados de servo con la calibración de cada articulación y
envía las doce órdenes a la interfaz de E/S (modo de transmisión continua del
firmware), que responde con todas las lecturas de sensores.

Cada articulación usa una cadena de fuentes de realimentación: su sensor
principal (potenciómetro leído a través del multiplexor analógico, o encoder
magnético AS5600); si falla, su sensor de respaldo; y si tampoco hay, la orden
retrasada 99~ms como estimación, Tabla~\ref{tab:sensores-articulaciones}. Las
cámaras, que identifican cada pieza de la pata por el color de su cinta (fémur
rojo, tibia verde, pie amarillo), se usan como tercera fuente, aunque solo en el
análisis posterior.

\begin{table}[H]
\centering
\caption{Fuentes de realimentación por articulación según la revisión del 28 de octubre de 2026.}
\label{tab:sensores-articulaciones}
\footnotesize
\begin{tabular}{llll}
\toprule
Articulación & Principal & Respaldo & Estado \\
\midrule
Fémur 1 (delantera derecha) & potenciómetro C8 & encoder ENC\_F1 & operativo \\
Fémur 2 (delantera izquierda) & potenciómetro C5 & encoder en 0x70/4 & operativo \\
Fémur 3 (trasera derecha) & potenciómetro C9 & --- & operativo \\
Fémur 4 (trasera izquierda) & potenciómetro C6 & encoder ENC\_F4 & operativo \\
Tibia 1 & --- & --- & sin sensor: orden retrasada \\
Tibia 2 & encoder ENC\_T2 (0x70/3) & --- & operativo; escala corregida \\
Tibia 3 & potenciómetro C12 & --- & responde, sin transmitir: orden retrasada \\
Tibia 4 & potenciómetro C10 & --- & operativo hasta el 5 de noviembre \\
\bottomrule
\end{tabular}
\end{table}

Las coxas se miden con los potenciómetros C0 a C3; el de la coxa 2 (C1) no sigue a
la coxa, por lo que su lectura no se considera fiable, y por eso el firmware no
corregía su caída.

\subsection{Correcciones necesarias para que el robot reproduzca la simulación}

Comparar el robot con la simulación mediante las cámaras reveló diferencias que
ningún entrenamiento podía compensar y que se corrigieron en la interfaz
robot--modelo:
\begin{itemize}
    \item \textbf{Signo de las tibias izquierdas.} Las tibias 2 y 4 están
    montadas en espejo: al pedirles levantar el pie lo bajaban. Se invirtió su
    signo; afectaba a todas las caminatas anteriores al 21 de octubre.
    \item \textbf{Postura base.} En la postura de reposo, la simulación tenía las
    rodillas dobladas unos $48^\circ$ y el robot casi rectas. Un ajuste
    automático en bucle, que mide con las cámaras el eje de la cinta de cada
    fémur y cada tibia y corrige el 70~\% de la diferencia por vuelta, dejó los
    fémures a menos de $2^\circ$ de la simulación y las tibias a unos $10^\circ$
    (antes, hasta $50^\circ$).
    \item \textbf{Topes del firmware.} Las correcciones de postura llevaban a
    las tibias 2 y 4 más allá de sus topes de seguridad, donde quedaban fijas.
    El programa respeta ahora esos topes y reserva margen para levantar el pie.
    \item \textbf{Escala de los sensores.} En la postura nueva los
    potenciómetros daban $3{,}2$ a $3{,}7$ cuentas por grado frente a las
    $2{,}5$ de la calibración anterior; la cámara confirmó que el motor sí se
    movía lo ordenado (ganancia $\approx1{,}08$), de modo que se recalibró la
    escala de los sensores y se mantuvo la ganancia del actuador en 1.
    \item \textbf{Temporización.} Con el PC cargado por el análisis de video y
    los entrenamientos, hasta 714 pasos de una caminata de 3~min llegaban con más
    de 40~ms de retraso. Se reservaron dos núcleos para el bucle de control, la
    red se evalúa con un único hilo y los análisis corren con prioridad mínima;
    con el PC entrenando a la vez, el número de pasos tardíos bajó a cero.
    \item \textbf{Postura de trabajo} (27 de octubre). El operador aprobó una postura
    nueva, con otras posiciones de apoyo y el movimiento de las tibias 2 y 4 invertido
    (sección de marchas nominales). Las comparaciones anteriores se repitieron por
    ello (sección~\ref{sec:postura27}). Ese día se detectó también que las políticas
    de gateo v14 y trote v15 se habían ejecutado sin la apertura de los fémures, porque
    un programa de selección dañó sus metadatos; se corrigió y se repararon los archivos.
    \item \textbf{Coxa 2} (28 de octubre). Se caía hacia adentro al caminar. Se ajustó
    en vivo con el operador hasta que dejó de caer: su base pasó de $48^\circ$ a
    $42^\circ$ en el firmware y en el programa de control.
    \item \textbf{Escala de la tibia 2} (28 de octubre). La escala de su encoder
    duplicaba la amplitud medida. Una regresión de la lectura frente a la orden en tres
    caminatas ($r\approx0{,}97$) la cambió de $-0{,}59$ a $-1{,}15$.
\end{itemize}

\section{Ciclo de mejora en la plataforma física}

Para que el robot no solo evaluara las políticas sino que también mejorara la
simulación, se implementó un ciclo automático (\path{ciclo_mejora.py}) con
cuatro pasos por vuelta:
\begin{enumerate}
    \item \textbf{Prueba en el robot}: la marcha nominal y las políticas
    candidatas de gateo y trote, dos repeticiones de 3~min cada una, en orden
    alternado entre rondas y con pausas de 2~min para enfriar los servos.
    \item \textbf{Puntuación con los sensores}: para cada condición se promedia
    el error RMS de seguimiento de fémures y tibias y la desviación media de las
    coxas, en grados. Gana la de menor puntuación.
    \item \textbf{Actualización del modelo}: con las caminatas de la vuelta se
    vuelven a medir, por articulación, el retardo, la histéresis y el error que
    el modelo no explica (\path{actualizar_modelo_real.py}). Se genera un modelo
    del actuador nuevo, con el intervalo de retardo medido para la
    aleatorización.
    \item \textbf{Reentrenamiento}: la política ganadora de cada marcha se sigue
    entrenando durante $10^6$ pasos con ese modelo, y su mejor copia se
    convierte en candidata de la vuelta siguiente. Se previeron cuatro semillas
    por marcha (401, 409, 419 y 421), pero por el error descrito en la
    Tabla~\ref{tab:campanas-ppo} las cuatro dieron la misma ejecución, de modo
    que en las vueltas realizadas cada marcha se reentrenó, en la práctica, con
    una sola semilla.
\end{enumerate}
De este modo la selección se decide con el robot real y no con la simulación, y
las discrepancias medidas en el robot se trasladan al entrenamiento. Una
limitación de este método es que la puntuación no mide la locomoción
propiamente dicha: el robot está suspendido, por lo que se evalúan el
seguimiento de las órdenes y la estabilidad de las coxas, no el avance ni la
inclinación del cuerpo sobre el suelo.
